\documentclass[utf8]{frontiers_suppmat}
\usepackage{url,hyperref,lineno,microtype,amsmath,booktabs}
\usepackage[onehalfspacing]{setspace}

\begin{document}
\onecolumn
\firstpage{1}

\title[Supplementary Material]{{\helveticaitalic{Supplementary Material: SIMPLEX Composition-Space Deep Learning for Hydrogel Adhesion}}}

\maketitle

\section{Supplementary Data}

This supplement provides the complete experimental and analytical details
behind the SIMPLEX framework summarised in the main article, including
dataset description, architecture configuration, training protocol, the full
ablation table, external ranking statistics, multi-target extensibility,
interpretation highlights, and a reproduction guide. All code and processed
data are available at \url{https://github.com/SHENTongfei/HydroGelNet}.

\section{Data and Preprocessing}

\subsection{Raw dataset}
The raw experimental data were obtained from the public repository
\emph{sheng-hu/hydrogels} (MIT licence) accompanying Liao et al.,
\textit{Nature} 644, 89--95 (2025), DOI 10.1038/s41586-025-09269-4. The
dataset records underwater glass adhesion strength (kPa) for 341 bio-inspired
hydrogel formulations, each described by the molar fractions of six
functional monomers:
\begin{itemize}
\item 2-hydroxyethyl acrylate (HEA, nucleophilic);
\item butyl acrylate (BA, hydrophobic);
\item 2-carboxyethyl acrylate (CBEA, acidic);
\item (3-acrylamidopropyl)trimethylammonium chloride (ATAC, cationic);
\item 2-phenoxyethyl acrylate (PEA, aromatic);
\item acrylamide (AAm, amide).
\end{itemize}
The six fractions sum to 1 (composition simplex).

\subsection{Internal / external split (model-guided extrapolation)}
\begin{itemize}
\item \textbf{Internal (training region):} the 180 round-1 bio-inspired
  formulations ($df\_180$). Adhesion mean 46.9 kPa, max 146.6 kPa.
\item \textbf{External (extrapolation region):} the 161 formulations added
  in later iterations of the Nature study's sequential model-based
  optimisation loop ($df\_341$ minus $df\_180$). Adhesion mean 154.2 kPa,
  max 353.3 kPa; 45\% of external samples exceed the training maximum.
\end{itemize}
This is a \emph{model-guided migration to a high-performance composition
region} (target-value extrapolation), not an independent-laboratory cohort.
The selection bias and range restriction inherent to the SMBO protocol are
acknowledged; range-restricted Spearman $\rho$ is a conservative estimate of
ranking ability.

\subsection{Features}
Two modalities, 21 features total: (1) \emph{monomer fractions} (6 raw molar
fractions); (2) \emph{pairwise synergy} (15 products $x_i x_j$, $i<j$).
Scaling/imputation were fitted inside each cross-validation training fold
only (no leakage).

\section{Model Configuration (final, best\_config.json)}

\begin{verbatim}
d_model=64, n_blocks=2, n_heads=4, dropout=0.2
use_transformer=False (ResBlock base), use_attention=True
use_film=False, use_modality_gate=False
use_mixup=True (alpha=0.4), use_swa=True
use_sam=False, use_ema=False
use_contrastive=False, use_pretrain_recon=False, use_mfm=False
use_uncertainty_weighting=False (single target)
use_domain_constraint=True (constraint_w=0.1)
y_transform=standard, scaler=standard
max_epochs=150, patience=30, batch_size=32
lr=3e-3, weight_decay=1e-3
\end{verbatim}

Design notes: (i) \textbf{ResBlock base, not Transformer} --- on 180 samples
the Transformer diverged (internal $R^2 \approx -6$); the residual network
with interaction attention is the stable configuration. (ii) \textbf{Mixup
is the key regulariser} --- removing it drops internal $R^2$ from 0.713 to
0.646. (iii) \textbf{Self-supervised pre-training (MFM/contrastive) was
detrimental} ($R^2 \approx -2.7$) and was removed (ablation-gated).

\section{Full Ablation Table (internal R2, 2 seeds x 3 folds)}

\begin{table}[h]
\centering
\begin{tabular}{lc}
\toprule
Variant & $\Delta R^2$ (full $-$ variant) \\
\midrule
w/o multimodal fusion & +0.093 (largest) \\
w/o Mixup & +0.067 \\
fusion = gated & +0.062 \\
w/o residual blocks & +0.062 \\
fusion = film & +0.031 \\
w/o task-specific gating & +0.022 \\
w/o attention & +0.019 \\
w/o domain constraint & +0.000 (neutral) \\
w/o SWA & +0.008 \\
\bottomrule
\end{tabular}
\end{table}

Positive delta means removal degrades performance (component helps). Sixteen
components that did not pay for themselves were removed from the final model;
the remaining switches are reported transparently.

\section{External Ranking Statistics}

Sample-level bootstrap (n = 161, 2000 resamples):
\begin{table}[h]
\centering
\begin{tabular}{lcc}
\toprule
Model & Spearman $\rho$ & 95\% CI \\
\midrule
SIMPLEX & 0.501 & [0.369, 0.619] \\
ElasticNet & 0.494 & [0.36, 0.62] \\
Ridge & 0.486 & [0.35, 0.61] \\
SVR-RBF & 0.379 & [0.24, 0.51] \\
MLP & 0.315 & [0.18, 0.45] \\
RandomForest & 0.211 & [0.06, 0.36] \\
\bottomrule
\end{tabular}
\end{table}

Paired bootstrap (SIMPLEX $-$ baseline):
\begin{itemize}
\item vs RandomForest: $\Delta\rho = +0.183$, 95\% CI $[+0.066, +0.311]$,
  $P(\Delta>0) = 0.999$ --- \textbf{significant};
\item vs ElasticNet: $\Delta\rho = -0.028$, 95\% CI $[-0.157, +0.097]$,
  $P(\Delta>0) = 0.345$ --- \textbf{statistical tie} (reported honestly).
\end{itemize}

Top-k screening precision (fraction of predicted top-k that are true top-k):

\begin{table}[h]
\centering
\begin{tabular}{lccc}
\toprule
$k$ & SIMPLEX & RF & ElasticNet \\
\midrule
20 & \textbf{0.25} & 0.10 & 0.05 \\
30 & \textbf{0.37} & 0.07 & 0.17 \\
\bottomrule
\end{tabular}
\end{table}

Prediction-range behaviour on the external set (kPa): SIMPLEX $[16, 109]$;
RandomForest $[26, 109]$ (compressed --- tree models are locked near the
training maximum $\approx 147$); ElasticNet $[-82, 232]$ (over-extrapolating,
non-physical negative predictions); true range $[3, 353]$. SIMPLEX is the
only model whose external predictions stay physically plausible
(non-negative, in-range) while retaining ranking signal.

\section{Multi-target Extensibility (preliminary)}

The original dataset also records rheological targets for the 180 internal
formulations: storage-related modulus (kPa) and loss tangent tan-delta.
$\mathrm{corr}(G, \mathrm{Modulus}) = -0.036$ --- independent information not
used by the Nature study's single-target models. A single-task SIMPLEX
trained on Modulus achieves internal CV $R^2 = 0.396$ vs RandomForest 0.392
(3 seeds), i.e. parity with the strongest baseline on a second, independent
target. Joint multi-target training is left for future work (see
Limitations).

\section{Interpretation Highlights}

Cross-validated permutation importance (mean over folds/seeds):
\begin{table}[h]
\centering
\begin{tabular}{lcc}
\toprule
Rank & Feature & Importance $\pm$ SE \\
\midrule
1 & Cationic-ATAC & 0.222 $\pm$ 0.028 \\
2 & pair (ATAC$\times$PEA) & 0.099 $\pm$ 0.017 \\
3 & pair (BA$\times$CBEA) & 0.072 $\pm$ 0.010 \\
4 & Hydrophobic-BA & 0.055 $\pm$ 0.013 \\
5 & pair (BA$\times$PEA) & 0.050 $\pm$ 0.020 \\
\bottomrule
\end{tabular}
\end{table}

The cation monomer and pairwise interaction terms dominate --- consistent
with the electrostatic/hydrophobic synergy known to drive underwater
adhesion, and with the SMBO-discovered high-performance region enriched in
the hydrophobic--aromatic (BA--PEA) combination.

\section{Reproducibility}

Environment: Python 3.11 (conda env \texttt{HydroGelNet}), PyTorch
2.14.0+cu130, NVIDIA RTX 5080 16 GB. Fixed seeds: [42, 2024, 7, 1337,
20260731]; all preprocessing inside folds. Full pipeline:
\texttt{download $\rightarrow$ build $\rightarrow$ qc $\rightarrow$ tune
$\rightarrow$ train $\rightarrow$ baselines $\rightarrow$ stats
$\rightarrow$ gate $\rightarrow$ interpret $\rightarrow$ figures
$\rightarrow$ tables $\rightarrow$ paper} (run\_all.py). PERF-GATE verdict:
PASS (G1 internal tie, G2 no significant disadvantage, G3 seed stability, G4
external Spearman 0.501 vs best baseline 0.494, G5 informative ablation).

\section{Data and Code Availability}

Data: \emph{sheng-hu/hydrogels} (MIT), Nature 2025 supplement. Code:
\url{https://github.com/SHENTongfei/HydroGelNet} (to be released).
Reproducible artefacts: \texttt{data/}, \texttt{results/},
\texttt{figures/}, \texttt{tables/}, \texttt{paper/}.

\end{document}
